\documentclass[10pt,a4paper]{amsart}
\usepackage{amsmath,amssymb,amsthm,graphicx}
% Retained minimal preamble for the Stacks Project chapter more-morphisms.tex (More on Morphisms).
% Verbatim excerpt of the chapter preamble: Xy-pic setup (native preamble lines 14-18) and the native
% theorem declarations with the mathematical macros (native preamble lines 154-217). Copied mechanically;
% nothing added, removed or redefined. The wrapper class and the amsmath/amssymb/amsthm/graphicx base
% packages are supplied by the assembler, so the class/externaldocument part of the native preamble is omitted.
% For commutative diagrams we use Xy-pic
\usepackage[all]{xy}

% We use 2cell for 2-commutative diagrams.
\xyoption{2cell}
% Theorem environments.
%
\theoremstyle{plain}
\newtheorem{theorem}[subsection]{Theorem}
\newtheorem{proposition}[subsection]{Proposition}
\newtheorem{lemma}[subsection]{Lemma}

\theoremstyle{definition}
\newtheorem{definition}[subsection]{Definition}
\newtheorem{example}[subsection]{Example}
\newtheorem{exercise}[subsection]{Exercise}
\newtheorem{situation}[subsection]{Situation}

\theoremstyle{remark}
\newtheorem{remark}[subsection]{Remark}
\newtheorem{remarks}[subsection]{Remarks}

\numberwithin{equation}{subsection}

% Macros
%
\def\lim{\mathop{\mathrm{lim}}\nolimits}
\def\colim{\mathop{\mathrm{colim}}\nolimits}
\def\Spec{\mathop{\mathrm{Spec}}}
\def\Hom{\mathop{\mathrm{Hom}}\nolimits}
\def\Ext{\mathop{\mathrm{Ext}}\nolimits}
\def\SheafHom{\mathop{\mathcal{H}\!\mathit{om}}\nolimits}
\def\SheafExt{\mathop{\mathcal{E}\!\mathit{xt}}\nolimits}
\def\Sch{\mathit{Sch}}
\def\Mor{\mathop{\mathrm{Mor}}\nolimits}
\def\Ob{\mathop{\mathrm{Ob}}\nolimits}
\def\Sh{\mathop{\mathit{Sh}}\nolimits}
\def\NL{\mathop{N\!L}\nolimits}
\def\CH{\mathop{\mathrm{CH}}\nolimits}
\def\proetale{{pro\text{-}\acute{e}tale}}
\def\etale{{\acute{e}tale}}
\def\QCoh{\mathit{QCoh}}
\def\Ker{\mathop{\mathrm{Ker}}}
\def\Im{\mathop{\mathrm{Im}}}
\def\Coker{\mathop{\mathrm{Coker}}}
\def\Coim{\mathop{\mathrm{Coim}}}

% Boxtimes
%
\DeclareMathSymbol{\boxtimes}{\mathbin}{AMSa}{"02}

%
% Macros for moduli stacks/spaces
%
\def\QCohstack{\mathcal{QC}\!\mathit{oh}}
\def\Cohstack{\mathcal{C}\!\mathit{oh}}
\def\Spacesstack{\mathcal{S}\!\mathit{paces}}
\def\Quotfunctor{\mathrm{Quot}}
\def\Hilbfunctor{\mathrm{Hilb}}
\def\Curvesstack{\mathcal{C}\!\mathit{urves}}
\def\Polarizedstack{\mathcal{P}\!\mathit{olarized}}
\def\Complexesstack{\mathcal{C}\!\mathit{omplexes}}
% \Pic is the operator that assigns to X its picard group, usage \Pic(X)
% \Picardstack_{X/B} denotes the Picard stack of X over B
% \Picardfunctor_{X/B} denotes the Picard functor of X over B
\def\Pic{\mathop{\mathrm{Pic}}\nolimits}
\def\Picardstack{\mathcal{P}\!\mathit{ic}}
\def\Picardfunctor{\mathrm{Pic}}
\def\Deformationcategory{\mathcal{D}\!\mathit{ef}}
\usepackage{cleveref}
\crefname{originalreference}{original reference}{original references}
\Crefname{originalreference}{Original reference}{Original references}
\makeatletter
\newlabel{more-algebra-lemma-relatively-pseudo-coherent}{{065H}{source}{}{sourcepointer-1}{}}
\newlabel{more-algebra-lemma-relatively-pseudo-coherent@cref}{{[originalreference][0][]065H}{[1][0][]source}{}{}{}}
\makeatother
\crefname{originalreference}{original reference}{original references}
\Crefname{originalreference}{Original reference}{Original references}
\makeatletter
\newlabel{perfect-lemma-closed-push-pseudo-coherent}{{09VA}{source}{}{sourcepointer-2}{}}
\newlabel{perfect-lemma-closed-push-pseudo-coherent@cref}{{[originalreference][0][]09VA}{[1][0][]source}{}{}{}}
\makeatother
\crefname{originalreference}{original reference}{original references}
\Crefname{originalreference}{Original reference}{Original references}
\makeatletter
\newlabel{perfect-lemma-pseudo-coherent-affine}{{08E7}{source}{}{sourcepointer-3}{}}
\newlabel{perfect-lemma-pseudo-coherent-affine@cref}{{[originalreference][0][]08E7}{[1][0][]source}{}{}{}}
\makeatother
\title{Relative pseudo-coherence tested after a closed immersion into affine space}
\author{The Stacks Project Authors}
\date{stacks-project:more-morphisms.tex:local-snapshot-20261011; GNU FDL 1.2 or later}
\begin{document}
\maketitle
\noindent\textit{Editorial note.} Counted span is the author's own proof environment, source lines 16625-16663, immediately following the statement at 16609-16623.\par\medskip
\noindent\textit{Editorial note.} Three cross-chapter pointers (More on Algebra, Derived Categories of Schemes) are recorded in assembly.external\_references.\par\medskip
\noindent\textit{Editorial note.} The statement's two conditions (some embedding / all embeddings) are kept verbatim in the statement span.\par\medskip
\noindent\textit{Allowed standard prerequisites (editorial).}\begin{itemize}
\item Ring-level criterion for relative pseudo-coherence. More on Algebra, lemma more-algebra-lemma-relatively-pseudo-coherent, cited in the proof. Supplies the algebraic counterpart of the equivalence, phrased for a complex of A-modules over R[x\_1..x\_n].
\item Pseudo-coherence and pushforward along a closed immersion. Derived Categories of Schemes, lemma perfect-lemma-closed-push-pseudo-coherent, cited in the proof. Moves pseudo-coherence across a closed immersion, which is what makes the choice of embedding irrelevant.
\item Affine comparison for pseudo-coherent objects. Derived Categories of Schemes, lemma perfect-lemma-pseudo-coherent-affine, cited in the proof. Identifies pseudo-coherence of the pushed-forward object with the module-level condition on global sections.
\end{itemize}\medskip
\ifdefined\hypertarget\hypertarget{sourcepointer-1}{}\fi
\noindent\textit{Source reference.} Original reference 065H: focused-seeds/geometry/S9-morph-relatively-pseudo-coherent-embedding/ref-tags.txt, source lines 4--4. Navigation pointer only: the counted proof cites this statement by \textbackslash{}ref as a standing input; the deduction the counted proof has to make is carried out inside the counted span, and no part of that deduction is delegated to this pointer.\par\medskip
\ifdefined\hypertarget\hypertarget{sourcepointer-2}{}\fi
\noindent\textit{Source reference.} Original reference 09VA: focused-seeds/geometry/S9-morph-relatively-pseudo-coherent-embedding/ref-tags.txt, source lines 5--5. Navigation pointer only: the counted proof cites this statement by \textbackslash{}ref as a standing input; the deduction the counted proof has to make is carried out inside the counted span, and no part of that deduction is delegated to this pointer.\par\medskip
\ifdefined\hypertarget\hypertarget{sourcepointer-3}{}\fi
\noindent\textit{Source reference.} Original reference 08E7: focused-seeds/geometry/S9-morph-relatively-pseudo-coherent-embedding/ref-tags.txt, source lines 6--6. Navigation pointer only: the counted proof cites this statement by \textbackslash{}ref as a standing input; the deduction the counted proof has to make is carried out inside the counted span, and no part of that deduction is delegated to this pointer.\par\medskip
\begin{lemma}
\label{lemma-relatively-pseudo-coherent}
Let $X \to S$ be a finite type morphism of affine schemes.
Let $E$ be an object of $D(\mathcal{O}_X)$.
Let $m \in \mathbf{Z}$.
The following are equivalent
\begin{enumerate}
\item for some closed immersion $i : X \to \mathbf{A}^n_S$
the object $Ri_*E$ of $D(\mathcal{O}_{\mathbf{A}^n_S})$
is $m$-pseudo-coherent, and
\item for all closed immersions $i : X \to \mathbf{A}^n_S$
the object $Ri_*E$ of $D(\mathcal{O}_{\mathbf{A}^n_S})$
is $m$-pseudo-coherent.
\end{enumerate}
\end{lemma}

\begin{proof}
Say $S = \Spec(R)$ and $X = \Spec(A)$. Let $i$ correspond to the surjection
$\alpha : R[x_1, \ldots, x_n] \to A$ and let $X \to \mathbf{A}^m_S$
correspond to $\beta : R[y_1, \ldots, y_m] \to A$.
Choose $f_j \in R[x_1, \ldots, x_n]$ with $\alpha(f_j) = \beta(y_j)$
and $g_i \in R[y_1, \ldots, y_m]$ with $\beta(g_i) = \alpha(x_i)$.
Then we get a commutative diagram
$$
\xymatrix{
R[x_1, \ldots, x_n, y_1, \ldots, y_m]
\ar[d]^{x_i \mapsto g_i} \ar[rr]_-{y_j \mapsto f_j} & &
R[x_1, \ldots, x_n] \ar[d] \\
R[y_1, \ldots, y_m] \ar[rr] & & A
}
$$
corresponding to the commutative diagram of closed immersions
$$
\xymatrix{
\mathbf{A}^{n + m}_S & \mathbf{A}^n_S \ar[l] \\
\mathbf{A}^m_S \ar[u] & X \ar[u] \ar[l]
}
$$
Thus it suffices to show that under a closed immersion
$$
f : \mathbf{A}^m_S \to \mathbf{A}^{n + m}_S
$$
an object $E$ of $D(\mathcal{O}_{\mathbf{A}^m_S})$ is
$m$-pseudo-coherent if and only if $Rf_*E$ is $m$-pseudo-coherent.
This follows from
Derived Categories of Schemes, Lemma
\ref{perfect-lemma-closed-push-pseudo-coherent}
and the fact that $f_*\mathcal{O}_{\mathbf{A}^m_S}$ is
a pseudo-coherent $\mathcal{O}_{\mathbf{A}^{n + m}_S}$-module.
The pseudo-coherence of $f_*\mathcal{O}_{\mathbf{A}^m_S}$ is
straightforward to prove directly, but it also follows from
Derived Categories of Schemes, Lemma \ref{perfect-lemma-pseudo-coherent-affine}
and
More on Algebra, Lemma \ref{more-algebra-lemma-relatively-pseudo-coherent}.
\end{proof}

\end{document}
